% ===========================================================================
% Poste 7 — Allocation famille + supplément fournitures scolaires
% ===========================================================================
\poste{7}{Allocation famille et supplément pour fournitures scolaires}{QC\_sae, SFS}

\pdesc Prestation versée par Retraite Québec pour chaque enfant à charge de
moins de 18~ans, bonifiée d'un supplément pour les familles monoparentales.
S'y ajoute, pour chaque enfant d'âge scolaire, un supplément annuel fixe
pour l'achat de fournitures scolaires.

\pobj Aider les familles à assumer les coûts liés aux enfants, avec un
soutien plancher universel : contrairement à plusieurs programmes, une
famille à revenu élevé reçoit quand même le montant minimal.

\pregle L'allocation part d'un montant maximal --- $M$ par enfant, plus le
supplément monoparental $S_{max}$ si le ménage ne compte qu'un adulte --- et
le réduit de \pc{4} du revenu familial net excédant le seuil $s$ ; mais elle
ne descend jamais sous le montant \emph{minimal} ($m$ par enfant, plus
$S_{min}$ si monoparental) :
\begin{align*}
\text{maximum} &= n \times M + \delta \times S_{max} \\
\text{minimum} &= n \times m + \delta \times S_{min} \\
\text{allocation} &= \max\bigl(\text{minimum},\;
\text{maximum} - 0{,}04 \times \pos{\RFN - s}\bigr)
\end{align*}
où $\delta = 1$ pour une famille monoparentale et $\delta = 0$ pour un
couple. Le seuil $s$ dépend de la composition (monoparentale ou couple). Le
montant par enfant ne dépend pas de son rang. Le résultat est arrondi au
cent.

\emph{Supplément pour fournitures scolaires.} Montant fixe versé pour chaque
enfant \textbf{de 4 à 16~ans} (âge $> 3$ et $< 17$), sans aucune réduction
selon le revenu :
\[
SFS = n_{4\text{--}16} \times F
\]

\pparams

\begin{center}
\small
\begin{tabular}{l r r}
\toprule
Paramètre & 2025 & 2026 \\
\midrule
Montant maximal par enfant ($M$)        & \D{3006} & \D{3068} \\
Montant minimal par enfant ($m$)        & \D{1196} & \D{1221} \\
Supplément monoparental maximal ($S_{max}$) & \D{1055} & \D{1077} \\
Supplément monoparental minimal ($S_{min}$) & \D{421}  & \D{430} \\
Seuil de réduction --- monoparentale ($s$)  & \D{43280} & \D{44032} \\
Seuil de réduction --- couple ($s$)         & \D{59369} & \D{60398} \\
Taux de réduction                       & \pc{4}   & \pc{4} \\
Fournitures scolaires, par enfant ($F$) & \D{124}  & \D{127} \\
\bottomrule
\end{tabular}
\end{center}

% ============================================================================
% GÉNÉRÉ par scripts/exemples-guide.ts — NE PAS ÉDITER À LA MAIN.
% Régénérer : npm run exemples (chaque chiffre est vérifié contre le moteur).
% ============================================================================
\pexemples Le montant maximal (par enfant, plus le supplément monoparental le cas
échéant) est réduit de \pc{4} du $\RFN$ excédant le seuil,
sans jamais descendre sous le montant minimal. Le supplément pour fournitures
scolaires (\D{124} par enfant de 4 à 16~ans) s'ajoute
sans égard au revenu.

\medskip\noindent\textbf{M6 --- famille monoparentale, 35~ans, revenu de travail de \D{35000}, un enfant : 3~ans (garde subventionnée, \D{2000}).}
$\RFN$ de \num{33265}, sous le seuil monoparental de \num{43280} :
aucune réduction, l'allocation est maximale.
\begin{align*}
\text{maximum} &= \num{3006} + \num{1055} = \num{4061} \\
\text{allocation} &= \num{4061} - 0 = \num{4061}
\end{align*}
L'enfant a 3~ans : pas de supplément pour fournitures scolaires
(\D{0} — il vise les 4 à 16~ans).

\medskip\noindent\textbf{M8 --- couple, 38 et 36~ans, revenus de travail de \D{60000} et \D{40000}, deux enfants : 4~ans (garde non subventionnée, \D{13000}) et 8~ans (garde non subventionnée, \D{3000}).}
\begin{align*}
\text{maximum} &= 2 \times \num{3006} = \num{6012} \\
\text{réduction} &= \pos{\num{96230} - \num{59369}} \times \pc{4} = \num{1474.44} \\
\text{allocation} &= \max\bigl(\num{2392},\ \num{6012} - \num{1474.44}\bigr) = \num{4537.56}
\end{align*}
Le plancher (montant minimal, \num{2392}) ne mord pas encore :
allocation de \D{4537.56}, plus \D{248} de fournitures scolaires
($2 \times \num{124}$, enfants de 4 et 8~ans).

\medskip\noindent\textbf{M9 --- couple, 45 et 45~ans, revenus de travail de \D{120000} et \D{60000}, un enfant : 5~ans (garde non subventionnée, \D{15000}).}
$\RFN$ de \num{175521} : la réduction (\num{4646.08}) dépasse
l'écart entre maximum (\num{3006}) et minimum (\num{1196}) —
l'allocation est au \emph{plancher}.
\begin{align*}
\text{allocation} &= \max\bigl(\num{1196},\ \num{3006} - \num{4646.08}\bigr) = \num{1196}
\end{align*}
Toute famille admissible reçoit au moins le minimum, ici \D{1196},
plus \D{124} de fournitures scolaires.


\prefs Loi sur les impôts (RLRQ, c.~I-3), art.~1029.8.61.8 à
1029.8.61.60 ; Retraite Québec ; ministère des Finances du Québec,
\emph{Paramètres du régime d'imposition}.
